\begin{tikzpicture}[
    font=\small,
    >=Stealth,
    detector/.style={circle, draw, fill=white, inner sep=0pt, minimum size=3.8mm, font=\scriptsize},
    active-detector/.style={circle, draw, fill=black, inner sep=0pt, minimum size=3.8mm}
]

    % ==========================================
    % 1. Malla de la región de imagen (Grid 9x9 EXACTA Y ENTERA)
    % ==========================================
    \begin{scope}[shift={(0,0)}]
        % Fondo blanco para limpiar
        \fill[white] (-2.25,-2.28) rectangle (2.25,2.22);

        % Dibujar las 8 líneas verticales internas
        \foreach \x in {-1.75, -1.25, -0.75, -0.25, 0.25, 0.75, 1.25, 1.75} {
            \draw[gray!50, thin] (\x, -2.25) -- (\x, 2.25);
        }

        % Dibujar las 8 líneas horizontales internas
        \foreach \y in {-1.75, -1.25, -0.75, -0.25, 0.25, 0.75, 1.25, 1.75} {
            \draw[gray!50, thin] (-2.25, \y) -- (2.25, \y);
        }

        % Borde exterior grueso de la grilla 9x9
        \draw[gray!80, thick] (-2.25,-2.25) rectangle (2.25,2.25);

        % Etiqueta N en la esquina superior derecha
        \node[below left=1pt, font=\small\itshape] at (2.25,2.25) {N};

        % Cotas dx y dy en celdas superiores completas
        \draw[<->, thin] (-1.1, 1.75) -- node[right, font=\tiny] {$dy$} (-1.1, 2.25);
        \draw[<->, thin] (-1.75, 1.55) -- node[font=\tiny, below] {$dx$} (-1.25, 1.55);

        % Contenido en los centros de celdas enteras
        \node[font=\tiny] at (-2.0, -2.0) {1};
        \node[font=\tiny] at (-1.5, -2.0) {2};
        \fill (0.5, -2) circle (1pt);
        \fill (-0.5, -2.0) circle (1pt);
        \fill (0.0, -2.0) circle (1pt);
        \node[font=\small\itshape] at (2.0, -2.0) {j};

        % Segunda fila inferior (Y = -1.5)
        \node[font=\tiny] at (-2, -1.5) {$j+1$};
        \node[font=\tiny] at (2.0, -1.5) {$2j$};

        % Puntos suspensivos columna izquierda
        \fill (-2, 0.5) circle (1pt);
        \fill (-2.0, -0.5) circle (1pt);
        \fill (-2.0, 0.0) circle (1pt);

        % Puntos suspensivos columna derecha
        \fill (2, 0.5) circle (1pt);
        \fill (2.0, -0.5) circle (1pt);
        \fill (2.0, 0.0) circle (1pt);

        % Ejes cartesianos partiendo del CENTRO EXACTO de la celda central (0,0)
        \draw[->] (0,0) -- (0.4,0) node[below, font=\small\itshape] {$x$};
        \draw[->] (0,0) -- (0,0.4) node[left, font=\small\itshape] {$y$};

        % Vector r_j hacia el centro de la celda j (2.0, -2.0)
        \draw[->, thick] (0,0) -- (2.0, -2.0) node[midway, above right, font=\large\bfseries] {$\boldsymbol{r}_j$};

        % Etiqueta inferior corregida
        \node[below, font=\Large] at (0, -2.4) {región de interés};
    \end{scope}

    % ==========================================
    % 2. ANILLO CIRCULAR DE DETECTORES
    % ==========================================
    
    \draw[dashed, thin] (0,0) circle (3.8cm);

    % Detectores
    \node[detector] (d_nd) at (115:3.8cm) {};
    \node[above=2pt, font=\small\itshape] at (d_nd.north) {$N_s$};

    \node[detector] (d_1) at (90:3.8cm) {};
    \node[above=2pt, font=\small] at (d_1.north) {1};

    \node[detector] (d_2) at (65:3.8cm) {};
    \node[above=2pt, font=\small] at (d_2.north) {2};

    \node[detector] at (40:3.8cm) {};
    \node[detector] at (-10:3.8cm) {};
    \node[detector] at (-35:3.8cm) {};
    \node[detector] at (-60:3.8cm) {};
    \node[detector] at (-85:3.8cm) {};
    \node[detector] at (-110:3.8cm) {};
    \node[detector] at (-135:3.8cm) {};
    \node[detector] at (-160:3.8cm) {};
    \node[detector] at (180:3.8cm) {};
    \node[detector] at (155:3.8cm) {};
    \node[detector] at (135:3.8cm) {};

    % Detector activo 'l' (negro)
    \node[active-detector] (d_l) at (20:3.8cm) {};
    \node[right=3pt, font=\large\itshape] at (d_l.east) {$l$};

    % Vector r_d
    \draw[->, thick] (0,0) -- (d_l.center) node[above left, font=\large\bfseries, very near end] {$\boldsymbol{r}_d$};

    % Puntos suspensivos
    \fill (28:4.1cm) circle (1.2pt);
    \fill (32:4.1cm) circle (1.2pt);
    \fill (36:4.1cm) circle (0.04cm);

    \fill (12:4.1cm) circle (1.2pt);
    \fill (8:4.1cm) circle (0.04cm);
    \fill (4:4.1cm) circle (0.04cm);

    % Etiqueta Sensores
    \draw[<-, line width=1pt] ([xshift=2pt, yshift=-2pt]d_nd.south east) -- ++(0.36,-0.3) 
        node[right, align=left, font=\Large] {Sensores};

    % ==========================================
    % 3. GRÁFICO SEÑAL FOTOACÚSTICA (p_d vs T)
    % ==========================================
    \begin{scope}[shift={($(d_l.east)+(0.6,0)$)}]
        % Ejes cartesianos
        \draw[->, thin] (-0.1,0) -- (2.2,0) node[below, font=\small] {$T$};
        \draw[->, thin] (0,-0.8) -- (0,1.2) node[left, font=\small] {$p_d$};

        % Curva del pulso ultrasónico
        \draw[blue!40!black, thick] (0,0) 
            -- (0.2,0)
            .. controls (0.35,0.2) .. (0.45,0.9)
            .. controls (0.55,-0.8) .. (0.65,0.3)
            .. controls (0.75,-0.3) .. (0.85,0.15)
            .. controls (0.95,-0.1) .. (1.05,0.05)
            -- (2.0,0);
    \end{scope}

\end{tikzpicture}
